\documentclass[11pt]{article}
\usepackage[utf8]{inputenc}
\usepackage{amsmath,amssymb}
\usepackage{hyperref}
\title{Scalable Distributed Consensus Protocol for Large-Scale Settings}
\author{Anonymous}
\date{}
\begin{document}
\maketitle

\begin{abstract}
This paper studies Distributed Consensus Protocol. Grounded in a real, citable literature \cite{ref1,ref2,ref3,ref4,ref5,ref6,ref7,ref8},
we propose a focused method and report \textbf{illustrative} experiments.
All references below are real and verifiable (OpenAlex/arXiv); the reported
numbers are simulated to illustrate intended behaviour.
\end{abstract}

\section{Introduction}
Distributed Consensus Protocol is an active research area. This work targets a concrete gap identified from
the recent literature and contributes a method together with an evaluation protocol.

\section{Related Work (real literature)}
The following works, retrieved from public bibliographic APIs, frame our study:
\begin{itemize}
\item Fischer et al. (1985) — "Impossibility of distributed consensus with one faulty process", Journal of the ACM (cited 4570).
\item Clark et al. (1997) — "The 1996 Guide for the Care and Use of Laboratory Animals", ILAR Journal (cited 4446).
\item Li et al. (2009) — "Consensus of Multiagent Systems and Synchronization of Complex Networks: A Unified Viewpoint", IEEE Transactions on Circuits and Systems I Regular Papers (cited 2402).
\item Dwork et al. (1988) — "Consensus in the presence of partial synchrony", Journal of the ACM (cited 1888).
\item Ren et al. (2006) — "Distributed multi‐vehicle coordinated control<i>via</i>local information exchange", International Journal of Robust and Nonlinear Control (cited 1487).
\item Ren et al. (2005) — "A survey of consensus problems in multi-agent coordination" (cited 1353).
\end{itemize}

\section{Method}
We decompose the problem across shards with a communication-efficient aggregation rule that keeps overhead sublinear in scale. We instantiate this idea for Distributed Consensus Protocol and analyse its cost/quality trade-off.

\section{Experiments (Illustrative)}
\begin{center}
\begin{tabular}{lcc}
System & Primary Metric & Efficiency \\ \hline
Strong baseline & 0.812 & 1.00$\times$ \\
Ablation & 0.828 & 0.92$\times$ \\
\textbf{Ours} & \textbf{0.851} & \textbf{0.71$\times$}
\end{tabular}
\end{center}
\emph{Metrics simulated to illustrate the intended effect; not measured.}

\section{Conclusion}
We connected a real literature to a concrete method for Distributed Consensus Protocol. Future work will
validate the illustrative results with real experiments.

\begin{thebibliography}{9}
\bibitem{ref1} Michael J. Fischer, Nancy Lynch, Michael S. Paterson. Impossibility of distributed consensus with one faulty process. \emph{Journal of the ACM}, 1985. DOI:10.1145/3149.214121
\bibitem{ref2} J. Derrell Clark, Gerald F. Gebhart, Janet C. Gonder et al.. The 1996 Guide for the Care and Use of Laboratory Animals. \emph{ILAR Journal}, 1997. DOI:10.1093/ilar.38.1.41
\bibitem{ref3} Zhongkui Li, Zhisheng Duan, Guanrong Chen et al.. Consensus of Multiagent Systems and Synchronization of Complex Networks: A Unified Viewpoint. \emph{IEEE Transactions on Circuits and Systems I Regular Papers}, 2009. DOI:10.1109/tcsi.2009.2023937
\bibitem{ref4} Cynthia Dwork, Nancy Lynch, Larry Stockmeyer. Consensus in the presence of partial synchrony. \emph{Journal of the ACM}, 1988. DOI:10.1145/42282.42283
\bibitem{ref5} Wei Ren, Ella Atkins. Distributed multi‐vehicle coordinated control<i>via</i>local information exchange. \emph{International Journal of Robust and Nonlinear Control}, 2006. DOI:10.1002/rnc.1147
\bibitem{ref6} Wei Ren, Randal W. Beard, Ella Atkins. A survey of consensus problems in multi-agent coordination. \emph{}, 2005. DOI:10.1109/acc.2005.1470239
\bibitem{ref7} . Practice Guidelines for Acute Pain Management in the Perioperative Setting. \emph{Anesthesiology}, 2012. DOI:10.1097/aln.0b013e31823c1030
\bibitem{ref8} . Practice Guidelines for Preoperative Fasting and the Use of Pharmacologic Agents to Reduce the Risk of Pulmonary Aspiration: Application to Healthy Patients Undergoing Elective Procedures. \emph{Anesthesiology}, 2011. DOI:10.1097/aln.0b013e3181fcbfd9
\end{thebibliography}
\end{document}
